% !TeX root = ../papers/nc-p3.tex
\section{四中心パラメータからの可除指数と特異曲線の横断型}
\label{sec:central-cut-index-and-transverse-type}
\begin{theorem}\label{thm:central-cut-division-index}
標準正則窓の折り畳みが四つの中心四次パラメータ上有限自由で、全中心 $Z$ が整、
一般商が中心単純とする。$P$ を四パラメータ環、$d=[\operatorname{Frac}Z:\operatorname{Frac}P]$
とすると、$d=1,4,16,64$ のいずれかである。射影中心の関数体での零番隅の次数を $s$ とすれば
$d s^2=64$ であり、隅は可除環である。全ての始点で合成可能な非零経路元の積は非零となる。
有限辺なら $s=1$ は起きない。
\end{theorem}
\begin{proof}
経路長を4で割った次数では各隅の $P$ 自由階数は64である。
実際対応する経路長 $4m+j-i$ の次元は二項係数で、$m^3/6$ の係数が64である。
中心の一次数元で次数方向を自明化しても一般階数は変わらない。
各頂点の行列階数を $s_i$ とすると対角隅で $d s_i^2=64$、従って全て同じ $s$ である。
$X=\operatorname{Proj}Z$ で偏極 $H$ は四パラメータ射による $\mathcal O(1)$ の引き戻しであり
$H^3=d$ である（有限射で、一次数元があるので一般次数は $d$）。
厳密には $Z$ が次数1で生成されない場合も、四パラメータが基点無しのため
この $H$ の局所自明化を使える。

射影関数体でのBrauer指数を $r$ とする。一つの一般単純右加群を
$X$ 上の有限な代数層の連接加群 $M$ に延ばす。例えば有理基底の格子を稠密なアフィン開集合で
代数層により生成し、その連接延長を取り、さらに代数の作用で閉じればよい。
その各頂点成分の一般 $\mathcal O_X$ 階数は $sr$ である。これを次数付き窓加群に戻した
高い次数のHilbert級数は $p(t)/(1-t)^4$, $p(t)\in\mathbb Z[t,t^{-1}]$ である。
中心上有限なので高い捩りの全切断から有限生成の加群を取れ、
正則性とNoether性により有限生成頂点射影の有限長分解があるからである。
各 $n=4m+j$ の次元で $m^3/6$ の係数を比べると
\[
             64 p(1)=d s r,\qquad 64=d s^2 .
\]
従って $r/s=p(1)$ は正整数で、指数が次数を割るため $r=s$ である。
各頂点の一般右イデアルはそれぞれ一つの単純なので、異なる頂点間でも非零写像の合成は非零である。
$P$ 上自由なので一般商への写像は単射である。
$s=1$ なら滑らかな稠密開集合で一般単純は各頂点一次元であり、異なる射影中心指標を持つ。
実際中心単純な有限代数は稠密開集合でAzumayaである。有限点グラフの四組は有限個で、
次数倍率を除くと有限個の射影指標しか与えないため矛盾する。
\end{proof}

\begin{theorem}\label{thm:central-cut-transverse-a-one}
前節の中心型と全中心切断の条件で、節点条件を外し、全中心には正規性だけを課す。
普通の四次超曲面でも滑らかな二次曲面の四次分岐二重被覆でもよい。
射影中心の特異曲線があれば、一般超平面によるその全ての幾何学的横断特異点は
$A_1$ である。完備した曲面上の代数は
$\operatorname{End}_R(R^{\oplus8}\oplus I^{\oplus8})$ となり、
$R=k[[u,v,w]]/(uv-w^2)$, $I$ は非自明な直既約極大Cohen--Macaulay加群である。
各頂点は $R^2\oplus I^2$ に対応する。この主張は中心の特別な孤立特異点や
曲線同士の交点に追加の型を課さない。
\end{theorem}
\begin{proof}
$K$ を $|H|$ の関数体、$S/K$ を一般超平面とする。
普遍接続は $X$ 上の射影三次元空間束で、そこへ代数層を引き戻して超平面パラメータの
一般点に局所化したものが $S$ 上の代数である。$S$ は幾何学的に整かつ正規で、
特異点は高々有限、Gorensteinで双対化層は自明である。
実際Bertiniと有限中心射により滑らかな開集合では滑らかであり、特異集合の次元は高々1、
一般切断はその一般点だけに会う。局所超曲面であって正則性は余次元1で成立する。
中心上の代数は一般商で中心単純、滑らかな所でAzumayaである
（前節の冒頭と同じFrobenius・traceの議論）。先の定理によりここでは $s=4$ である。

曲面上の代数は有限大域次元を持つ。この点には特定の閉超平面での商の有限大域次元を使わない。
接続の全空間は射影空間束の各チャートで元の代数の多項式拡大の局所化だからである。
元の代数層の各中心アフィン開集合での有限大域次元は四次パラメータを可逆化して
次数を除く議論による。局所化、中心エタール変更、一般体 $K$ の分離代数拡大でも
この有限性は保たれる。幾何学的な一般切断には有限表示の分解と平坦移行を使えばよい。
代数はGorenstein中心上極大Cohen--Macaulayで、その中心双対は被約traceにより代数自身である。
これらの性質は平坦な接続・局所化で保たれる。局所の中心双対を使うと各閉点の単純 $L$ に対し
$\operatorname{Ext}^j_A(L,A)=\operatorname{Ext}^j_R(L,R)$ で、非零次数は2だけである。
有限射影次元と局所最小分解（または完備化）から射影次元は2となる。
従って幾何学的一般切断のアフィンチャートの代数はhomologically homogeneousである。
Stafford--Van den Berghの標数零の有理特異点定理をこの素な中心有限代数に適用すると、
中心曲面は有理特異点を持つ。素性は中心上のtorsion-free性と一般中心単純性から従う。
従ってGorenstein曲面の各特異点は有理二重点である。

その完備局所環 $R$ の穴あき部分上のBrauer群は零である。この局所事実も説明する。
最小解消の例外は滑らかな有理曲線の木である。穴あき部分のAzumaya類の分数体での
剰余は例外因子にしか存在しない。正則曲面の剰余複体で隣り合う剰余の合成は零なので、
木の末端の $\mathbb P^1$ 上の指標の分岐候補は隣の一交点だけである。
代数閉剰余体上で一か所以内で分岐する巡回被覆は自明である
（Kummer表示の因子の次数零からも従う）。末端を除いて繰り返すと全ての剰余は零となる。
純性で解消へ延び、固有基底変換でその $H^2(\mu_n)$ は例外各曲線の次数の直和である。
多重次数を指定した直線束は例外の厚みに延長できる（一次元の層の第二コホモロジーは零）。
形式存在定理により完備底上の解消に延びるので、Kummer列からこのBrauer類は零となる。

そこで分数体は分裂し、完備代数は $\operatorname{End}_R(M)$ と書ける。
穴あき部分で行列環の分裂束を取り、反射的に延長すればよい。
高さ1での一致と反射性により自己準同型環まで一致する。
任意の反射的加群 $N$ に対し $\operatorname{Hom}_R(M,N)$ は反射的で、曲面では
極大Cohen--Macaulayである。中心双対の式より、この加群から $A$ への正次数Extは
消える。有限大域次元のもと最小分解によりこの加群は射影的である。
よって $N\in\mathrm{add}(M)$ である。事実、射影的なHomは $M$ の有限直和の
ある直和因子 $N'$ のHomと同型であり、穴あき部分で森田同型を使い反射性で $N\simeq N'$ と戻せる。

有理二重点を $k[[x,y]]^G$, $G\subset\mathrm{SL}_2$ と表示する。
その反射的直既約加群は $G$ の各既約表現に対応し、階数はその次元である。
これは穴あき正則被覆に引き戻して反射的に延長すると自由になり、
剰余表現を平均化で自由生成元へ同変に持ち上げられることによる。
よって全既約表現が $M$ に現れる。曲面上のこの点の各単純も元の窓の基礎拡大後の単純である。
各頂点の次元は等しく、各単純でその数を $d_j$ とする。
有限の薄い四組の中心指標は一般超平面に入らないので全て $d_j\ge2$ である。
一頂点の行列階数は4であり、従って
$4=\sum_j d_j\dim\rho_j$ となる。非自明な群には少なくとも二既約表現があるので、
二つの一次元表現だけがあり、両 $d_j=2$ である。群の位数は次元平方和2、
従って $G=\{1,-1\}$、主張の横断型を得る。
\end{proof}
